\RequirePackage{pdfmanagement}
\SetKeys[document/metadata]{pdfversion=1.7, pdfstandard=A-2u, lang=en-GB}
\input{preamble/pdfA-compliance/glyphtounicode}
\documentclass[11pt,a4paper]{report}
\newcommand*{\PaperTitle}{The Kaibel--Petlyuk vapour penalty: a previously answered candidate}
\newcommand*{\PaperTitlePlain}{The Kaibel-Petlyuk vapour penalty: a previously answered candidate}
\newcommand*{\PaperShortTitle}{The Kaibel--Petlyuk vapour penalty}
\newcommand*{\PaperSubtitle}{chemistry-open-problems \textperiodcentered{} 27 October 2026 \textperiodcentered{} already solved}
\newcommand*{\PaperCredits}{Original result: Ivar J. Halvorsen and Sigurd Skogestad (2006).\\
Literature audit, exposition and certificates: GPT-6-Astra (effort High).}
\newcommand*{\PaperKeywords}{Kaibel column, Petlyuk arrangement, Underwood equations, minimum vapour flow}
\input{chemop-preamble}
\addbibresource{references.bib}

\begin{document}
\makepapertitle
\begin{paperabstract}
The repository shortlisted the strict energy penalty of the four-product Kaibel
column relative to the fully extended Petlyuk arrangement, together with an
analytic expression for that penalty. The literature check resolves this candidate:
Halvorsen and Skogestad answered it in 2006 in the ideal, infinite-stage Underwood
model. This note records their priority, makes the strict comparison explicit,
and supplies exact rational checks. We correct the printed minimum in their
equation (7) to a maximum, as required by their accompanying text and example.
This is an exposition of a known result, not a new solution of an open problem.
\end{paperabstract}

\section{The candidate and its disposition}
The candidate in the repository asked whether the Kaibel arrangement has a
strictly positive energy penalty for every four-component feed and whether that
penalty has a closed form. No source declaring this question open was attached.
The primary source located during the audit is Halvorsen--Skogestad
\autocite{HS2006}, Section 2, PDF page 4:
\begin{openquote}
The minimum vapour flow for the Kaibel arrangement can never come down to the
minimum energy for the extended four-product Petlyuk column.
\end{openquote}
Their Section 2 gives the algebraic calculation, and Section 4 compares the
relevant peaks. Thus the correct repository status is \emph{already solved}.
The exposition and certificates below document that disposition; they do not
claim priority for the comparison or the underlying column model.

\section{Model and exact answer}
Assume a positive feed rate \(F\), four strictly positive mole fractions
\(z_A,z_B,z_C,z_D\) summing to one, and constant relative volatilities
\[
 a>b>c>d>0.
\]
The products are pure, pressure and molar flows are constant in the sense of
the Underwood model, and the number of equilibrium stages tends to infinity.
Write \(q\) for the feed thermal parameter and \(s=1-q\). The algebra below
holds for every finite real \(q\); the usual two-phase range is \(0\le q\le1\).
We compare normalized \emph{overhead vapour flows}, not exergy or capital cost.
For a common feed the corresponding bottom flows subtract the same \(sF\),
so their difference is identical. Conversion to heat duties requires the
same latent-heat convention for both arrangements.

Define the two rational functions
\begin{equation}\label{eq:hg}
 h(x)=\frac{az_A}{a-x}+\frac{bz_B}{b-x},\qquad
 g(x)=\frac{cz_C}{c-x}+\frac{dz_D}{d-x}.
\end{equation}
Let \(\theta_A,t,\theta_C\) be the unique roots of \(h(x)+g(x)=s\) in
\((b,a),(c,b),(d,c)\), respectively. Put
\begin{equation}\label{eq:uv}
 u=h(t)>0,\qquad v=u-s=-g(t)>0.
\end{equation}
The three normalized Petlyuk peak requirements are
\begin{equation}\label{eq:peaks}
 H_A=\frac{az_A}{a-\theta_A},\qquad H_B=u,\qquad
 H_C=s+\frac{dz_D}{\theta_C-d}.
\end{equation}
The last expression also equals the sum of the A, B and C terms in the
feed equation evaluated at \(\theta_C\).

\begin{theorem}[note=Known comparison in the Underwood model]\label{thm:gap}
Under the stated assumptions the ideal minimum overhead flows are
\begin{align}
 P:=\frac{V_{T,\min}^{\rm Petlyuk}}F
     &=\max\{H_A,u,H_C\},\label{eq:P}\\
 K:=\frac{V_{T,\min}^{\rm Kaibel}}F
     &=\max\left\{\frac{u(a-t)}{a-b},\;
                 s+\frac{v(t-d)}{c-d}\right\}.\label{eq:K}
\end{align}
Consequently the penalty is the explicit algebraic quantity
\begin{equation}\label{eq:gap}
 \Delta V_T=F(K-P)>0.
\end{equation}
The same formula gives \(\Delta V_B\). The strict inequality is pointwise;
no feed-independent positive lower bound is asserted.
\end{theorem}

Here ``algebraic'' has a precise meaning: with
\(\alpha=(a,b,c,d)\), the three specified roots are roots of
\begin{equation}\label{eq:polynomial}
 R(x)=\sum_{i=1}^4\alpha_i z_i
             \prod_{j\ne i}(\alpha_j-x)
             -s\prod_{i=1}^4(\alpha_i-x).
\end{equation}
It has degree four when \(s\ne0\) and degree three when \(s=0\).
The intervals specify each required root uniquely, eliminating the ambiguity
of a bare polynomial-root expression. Thus only a polynomial of degree at
most four and rational arithmetic are needed; a choice of quartic or cubic
radicals can be substituted if desired. Formula \eqref{eq:gap} is not a
claim that the penalty is rational in the input data.

\section{Derivation and strictness}
\paragraph{Existence and location of the feed roots.}
Every term \(\alpha_i z_i/(\alpha_i-x)\) has positive derivative away
from its pole. On each interval between two consecutive volatilities,
\(h+g\) increases continuously from \(-\infty\) to \(+\infty\).
This proves the asserted existence and uniqueness for any real \(s\).
For \(c<t<b\), both terms of \(h(t)\) are positive and both terms of
\(g(t)\) are negative, proving \eqref{eq:uv}.

\paragraph{The column requirements used from the primary source.}
The ideal fully extended Petlyuk arrangement has requirement \eqref{eq:P}.
For the Kaibel arrangement the prefractionator makes the sharp AB/CD split,
with minimum top flow \(uF\) and bottom flow \(vF\).
The propagated roots are the unique numbers
\begin{equation}\label{eq:propagated}
 \phi\in(b,a),\quad h(\phi)=u,\qquad
 \psi\in(d,c),\quad g(\psi)=-v.
\end{equation}
The binary sections therefore require, in overhead-flow units,
\begin{equation}\label{eq:binary}
 K_A=\frac{az_A}{a-\phi},\qquad
 K_C=s+\frac{dz_D}{\psi-d}.
\end{equation}
These are the thermal-coupling and Underwood relations of
\autocite{HS2006}, equations (1)--(6). We use that ideal model as the
premise of the comparison, rather than asserting a new derivation of
distillation theory from stage equations.
Both sections must operate, so the total requirement is
\(K=\max\{K_A,K_C\}\). A minimum would satisfy only the less demanding
section. This corrects the two occurrences of \(\min\) in the source's
equation (7), consistently with its Sections 2 and 4 and Table 1.

Operating the prefractionator above its minimum cannot lower this optimum
within the same model. Indeed, replace \(u\) by \(w\ge u\) in
\(h(\phi_w)=w\) and \(g(\psi_w)=s-w\). Then \(\phi_w\) increases
and \(\psi_w\) decreases with \(w\), so both binary requirements in
\eqref{eq:binary} increase. Hence their maximum is minimized at \(w=u\).

\paragraph{Eliminating the propagated roots.}
After multiplication by \((a-x)(b-x)\), the equation \(h(x)=u\) becomes
\[
 u(a-x)(b-x)-az_A(b-x)-bz_B(a-x)=u(x-t)(x-\phi).
\]
Both sides are quadratic with leading coefficient \(u\) and the same
two distinct roots. Evaluating at \(x=a\) gives
\[
 u(a-t)(a-\phi)=az_A(a-b),\qquad
 K_A=\frac{u(a-t)}{a-b}.
\]
Similarly, clearing the denominators in \(g(x)=-v\) gives
\[
 v(c-x)(d-x)+cz_C(d-x)+dz_D(c-x)=v(x-t)(x-\psi).
\]
At \(x=d\), this yields
\[
 v(t-d)(\psi-d)=dz_D(c-d),\qquad
 K_C=s+\frac{v(t-d)}{c-d}.
\]
This proves \eqref{eq:K} from \eqref{eq:binary} without additional
quadratic solves.

\paragraph{Comparison of the peaks.}
Since \(g\) is strictly increasing on \((c,\infty)\) and
\(\theta_A>t\),
\[
 h(\theta_A)=s-g(\theta_A)<s-g(t)=u=h(\phi).
\]
Strict increase of \(h\) on \((b,a)\) implies
\(\phi>\theta_A\), hence \(K_A>H_A\).
Similarly \(h\) is strictly increasing on \((-\infty,b)\) and
\(\theta_C<t\), so
\[
 g(\theta_C)=s-h(\theta_C)>s-h(t)=-v=g(\psi).
\]
Thus \(\psi<\theta_C\) and \(K_C>H_C\).
Finally the reduced formulas give the particularly simple differences
\begin{equation}\label{eq:strict}
 K_A-u=\frac{u(b-t)}{a-b}>0,\qquad
 K_C-u=\frac{v(t-c)}{c-d}>0.
\end{equation}
Whichever of \(H_A,u,H_C\) attains \(P\), one of \(K_A,K_C\)
is strictly larger. This proves \(K>P\), including ties between any
of the Petlyuk peaks, and completes the comparison.

\section{Exact computational checks}
The proof above covers the whole stated parameter domain. The examples
are reproducibility checks, not a finite substitute for that proof.
For the source's Table 1 feed,
\[
 F=1,\quad \alpha=(6,4,2,1),\quad
 z=(3/10,1/5,1/5,3/10),\quad q=1,
\]
the certificates enclose \(P\) between \(1.375\) and \(1.385\)
and \(K\) between \(1.825\) and \(1.835\), reproducing the published
rounded values \(1.38\) and \(1.83\).
The committed certificate contains much narrower rational enclosures,
of width at most \(10^{-10}\), for all three Petlyuk peaks, both
Kaibel peaks, the two overall requirements and their difference.

There are seven cases: that source example, saturated vapour,
a two-phase feed, an example with the bottom section as bottleneck,
an exactly rational middle root, and thermal parameters on either
side of \([0,1]\). All use \(\alpha=(6,4,2,1)\).
The rational-middle-root case has the source composition and
\(q=3/20\), giving \(t=3\), \(u=7/5\), \(v=11/20\),
\(K_A=21/10\), and \(K_C=39/20\).

The generator uses rational bisection. The checker does not trust that
procedure: it reads the saved root intervals, checks their endpoint
signs and locations, then propagates intervals using only integer
arithmetic and Python's \texttt{Fraction}. It checks all strict comparisons,
the stated decimal-rounding windows, and both the propagated-root and
reduced expressions. Intervals containing a denominator zero are rejected.
It exits with a nonzero code and prints \texttt{FAIL} on a failed check.
Run from the repository root:
\begin{verbatim}
python problems/kaibel-petlyuk-energy-gap/scripts/verify.py
\end{verbatim}
No external Python package is required. Regeneration is optional:
\begin{verbatim}
python problems/kaibel-petlyuk-energy-gap/scripts/generate.py
\end{verbatim}

\section{Scope and verification status}
Strictly positive amounts of all four components and distinct volatilities
are essential to this statement. Removing a component or merging
volatilities changes the problem and invalidates the strict-inequality
argument. Perfectly pure products and minimum vapour flows here are
infinite-stage limits; this note claims no attainment at a finite tray count.
It also makes no universal claim for nonideal mixtures, variable molar
flows, pressure differences, or impure-product specifications.

The literature result is credited to Halvorsen and Skogestad (2006).
The present exposition and certificate checks were performed by
GPT-6-Astra (effort High), with attribution supplied by the user.
No separate fresh-context verification was performed. The repository
status \emph{already solved} records the identified prior answer and
does not label this note as a newly verified research result.

\paperreferences
\end{document}
